% Topic from _Linear Algebra_ Jim Hefferon
%  http://joshua.smcvt.edu/linearalgebra
%  2012-Feb-12
\topic{网页排名}
\index{page ranking|(}

设想你想找一本关于线性代数的最好的书。
你很可能会尝试使用诸如 Google 这样的万维网搜索引擎。
这类引擎按照重要性对网页进行排序。
正如 Google 的创始人所说（\cite{BrinPage}），排名的定义是：
一个网页是重要的，如果其他重要网页链接到它：
"如果有许多网页指向某个网页，或者有一些指向它的网页本身
具有很高的 PageRank，那么该网页就可以有很高的 PageRank。"
但这难道不是循环的么\Dash 在尚未确定哪些网页重要之前，
怎么能判断一个网页是否重要？
答案是使用特征值与特征向量。

我们将介绍网页排名算法的一个简化版本。
为此，我们将万维网建模为通过链接相互连接的网页的集合。
这幅取自 \cite{Wills} 的图把网页画成圆圈，把链接画成箭头。
网页~$p_1$ 有一条链接指向网页~$p_2$。
网页~$p_2$ 有一条链接指向~$p_3$。
而 $p_3$ 有链接指向 $p_1$、$p_2$ 和~$p_3$。
% \begin{center}  % add a little vertical spacing; looked tight to me.
%   \shortstack{\rule{0em}{1.5ex} \\ \includegraphics{jc/mp/ch5.9} \\ \rule{0em}{1.5ex}}
% \end{center}
\begin{center} 
  \includegraphics{jc/mp/ch5.9}
\end{center}

关键思想是：如果一个网页经常被其他网页引用，
那么它就应该获得较高的排名。
% (In practice people have tried to game the ranking system by
% setting up link farms of many pages that all point to each other.
% But our model will be simple.)
也就是说，如果网页~$p_i$ 被网页~$p_j$ 链接到，
我们就提高网页~$p_i$ 的重要度。
该增量等于进行链接的网页~$p_j$ 的重要度
除以该网页上出链接的数目 $a_j$。
\begin{equation*}
  \mathcal{I}(p_i)=\sum_{\text{链入的网页~$p_j$}}  \frac{\mathcal{I}(p_j)}{a_j}
\end{equation*}
于是，$p_1$ 的重要度等于 $1/3$ 乘以~$p_3$ 的重要度，
因为指向~$p_1$ 的唯一链接来自~$p_3$。
类似地，~$p_2$ 的重要度是~$p_1$ 的重要度
加上 $1/3$ 乘以~$p_3$ 的重要度之和。

这一信息储存于下面的矩阵之中。
\begin{equation*}
  \begin{mat}
    0   &0  &1/3  &0   \\
    1   &0  &1/3  &0   \\
    0   &1  &0    &0 \\
    0   &0  &1/3  &0
  \end{mat}
\end{equation*}
该算法的发明者给出了一种理解这个矩阵的方式。
\begin{quotation}
可以将 PageRank 看作用户行为的一个模型。
我们假设有一个"随机冲浪者"，他随机地得到一个网页，
然后不停地点击链接，从不点击"后退"\ldots{}
% but eventually gets bored
% and starts on another random page. 
随机冲浪者访问某个网页的概率就是该网页的 PageRank。
\cite{BrinPage}
\end{quotation}
于是，考察矩阵的第一行可以发现，随机冲浪者
到达 $p_1$ 的唯一途径是来自~$p_3$。
该网页有三条链接，所以点击进入 $p_1$ 的概率是~$1/3$
乘以位于网页~$p_3$ 上的概率。

这就引出了网页~$p_4$ 的问题。
在互联网上，许多网页是\definend{悬空链接}\index{link!dangling}\index{dangling link}
或\definend{汇点链接}，\index{link!sink}\index{sink link}
即没有任何出链接。
访问到这种网页的随机冲浪者会发生什么？
最简单的做法是设想冲浪者完全随机地选择下一个网页。
\begin{equation*}
  H=\begin{mat}
    0   &0  &1/3  &1/4   \\
    1   &0  &1/3  &1/4   \\
    0   &1  &0    &1/4 \\
    0   &0  &1/3  &1/4
  \end{mat}
\end{equation*}

我们将求出向量 $\vec{\mathcal{I}}$，其分量是每个网页的
重要度排名 $\mathcal{I}(p_i)$。
采用这一记号，
我们对网页排名的要求就是
$H\vec{\mathcal{I}}=\vec{\mathcal{I}}$。
也就是说，我们要求的是该矩阵对应于
特征值~$\lambda=1$ 的特征向量。

下面是 \textit{Sage} 对特征向量的计算
（为适合排版而作了删节）。
\begin{lstlisting}
sage: H=matrix([[0,0,1/3,1/4], [1,0,1/3,1/4], [0,1,0,1/4], [0,0,1/3,1/4]])   
sage: H.eigenvectors_right()                                                 
[(1, [(1, 2, 9/4, 1)], 1), 
 (0, [(0, 1, 3, -4)], 1), 
 (-0.3750000000000000? - 0.4389855730355308?*I, 
      [(1, -0.1250000000000000? + 1.316956719106593?*I, 
           -1.875000000000000? - 1.316956719106593?*I, 1)], 1), 
 (-0.3750000000000000? + 0.4389855730355308?*I, 
      [(1, -0.1250000000000000? - 1.316956719106593?*I, 
           -1.875000000000000? + 1.316956719106593?*I, 1)], 1)]
\end{lstlisting}
\textit{Sage} 给出的对应于
特征值~$\lambda=1$ 的特征向量是这一个。
\begin{equation*}
  \colvec{1 \\ 2 \\ 9/4 \\ 1}
\end{equation*}
当然，该特征空间中有许多向量。
为得到网页排名数值，我们将其归一化为长度一。
\begin{lstlisting}
sage: v=vector([1, 2, 9/4, 1])
sage: v/v.norm()
(4/177*sqrt(177), 8/177*sqrt(177), 3/59*sqrt(177), 4/177*sqrt(177))
sage: w=v/v.norm()
sage: w.n()
(0.300658411201132, 0.601316822402263, 0.676481425202546, 0.300658411201132)
\end{lstlisting}
于是我们把第一个和第四个网页评为
同等重要。
第二和第三个网页的排名比它们高得多，
而这两个网页彼此的重要性大致相当。

我们再添加一项改进。
我们允许冲浪者随机挑选一个新网页，
即使他们并不在悬空网页上。
在下式中，这种情况以概率~$1-\alpha$ 发生。
\begin{equation*}
  G=\alpha\cdot\begin{mat}
    0   &0  &1/3  &1/4   \\
    1   &0  &1/3  &1/4   \\
    0   &1  &0    &1/4 \\
    0   &0  &1/3  &1/4
  \end{mat}
  +(1-\alpha)\cdot\begin{mat}
    1/4   &1/4  &1/4  &1/4   \\
    1/4   &1/4  &1/4  &1/4   \\
    1/4   &1/4  &1/4  &1/4  \\
    1/4   &1/4  &1/4  &1/4
  \end{mat}
\end{equation*}
这就是\definend{Google 矩阵}\index{Google matrix}\index{matrix!Google}。

实践中 $\alpha$ 通常取 $0.85$ 到~$0.99$ 之间的值。
下面是四个网页在各种
$\alpha$ 下的排名。
\begin{center}
  \begin{tabular}{r|cccc}
    \multicolumn{1}{c}{$\alpha$}
          &$0.85$  &$0.90$  &$0.95$ &$0.99$ \\ \hline
    $p_1$ &$0.325$ &$0.317$ &$0.309$  &$0.302$  \\
    $p_2$ &$0.602$ &$0.602$ &$0.602$  &$0.601$  \\
    $p_3$ &$0.652$ &$0.661$ &$0.669$  &$0.675$  \\
    $p_4$ &$0.325$ &$0.317$ &$0.309$  &$0.302$  
  \end{tabular}
\end{center}
% Numbers computed as:
% sage: H=matrix(QQ,[[0,0,1/3,1/4], [1,0,1/3,1/4], [0,1,0,1/4], [0,0,1/3,1/4]])
% sage: S=matrix(QQ,[[1/4,1/4,1/4,1/4], [1/4,1/4,1/4,1/4], [1/4,1/4,1/4,1/4], [1/4,1/4,1/4,1/4]]) 
% sage: I=matrix(QQ,[[1,0,0,0], [0,1,0,0], [0,0,1,0], [0,0,0,1]])
% sage: alpha=0.95                       
% sage: G=alpha*H+(1-alpha)*S    
% sage: N=G-I  
% sage: 1200*N         
% [-1185.00000000000  15.0000000000000  395.000000000000  300.000000000000]
% [ 1155.00000000000 -1185.00000000000  395.000000000000  300.000000000000]
% [ 15.0000000000000  1155.00000000000 -1185.00000000000  300.000000000000]
% [ 15.0000000000000  15.0000000000000  395.000000000000 -900.000000000000]
% sage: M=matrix(QQ,[[-1185,15,395,300], [1155,-1185,395,300], [15,1155,-1185,300], [15,15,395,-900]])
% sage: M.echelon_form()         
% [         1          0          0         -1]
% [         0          1          0     -39/20]
% [         0          0          1 -3423/1580]
% [         0          0          0          0]
% sage: v=vector([1,39/20,3423/1580,1]) 
% sage: v/v.norm().n()   
% (0.308665053665036, 0.601896854646821, 0.668709163731278, 0.308665053665036)


\medskip
商业搜索引擎所用算法的细节是保密的，
它们有许多改进，而且变化频繁。
但 Google 的发明者非常大方地在 \cite{BrinPage} 中
概述了他们工作的基础。
较新的资料来源是 \cite{WikipediaPageRank}。
另外两篇出色的阐述是
\cite{Wills} 和
\cite{Austin}。

\begin{exercises}
  \item 若一个方阵的每一列元素之和为一，
    则称它是\definend{随机的}\index{stochastic matrix}\index{matrix!stochastic}。
    Google 矩阵由两个随机矩阵的组合
    $G=\alpha*H+(1-\alpha)*S$ 计算而来。
    证明 $G$ 必是随机的。
    \begin{answer}
      第~$j$ 列的元素之和为
      $\sum_i \alpha h_{i,j}+(1-\alpha)s_{i,j}
       =\sum_i \alpha h_{i,j}+\sum_i (1-\alpha)s_{i,j}
       =\alpha\sum_i \alpha h_{i,j} +(1-\alpha)\sum_i s_{i,j}
       =\alpha\cdot 1+(1-\alpha)\cdot 1$，
      其结果为一。
    \end{answer}
  \item 对于这个网页构成的网络，每个网页的重要度应当相等。
    请在 $\alpha=0.85$ 时验证这一点。
    \begin{center}
      \includegraphics{jc/mp/ch5.11}
    \end{center}
    \begin{answer}
      下面的 \textit{Sage} 会话给出了相等的值。
\begin{lstlisting}
sage: H=matrix(QQ,[[0,0,0,1], [1,0,0,0], [0,1,0,0], [0,0,1,0]])
sage: S=matrix(QQ,[[1/4,1/4,1/4,1/4], [1/4,1/4,1/4,1/4], [1/4,1/4,1/4,1/4], [1/4,1/4,1/4,1/4]])
sage: alpha=0.85     
sage: G=alpha*H+(1-alpha)*S    
sage: I=matrix(QQ,[[1,0,0,0], [0,1,0,0], [0,0,1,0], [0,0,0,1]])
sage: N=G-I
sage: 1200*N         
[-1155.00000000000  45.0000000000000  45.0000000000000  1065.00000000000]
[ 1065.00000000000 -1155.00000000000  45.0000000000000  45.0000000000000]
[ 45.0000000000000  1065.00000000000 -1155.00000000000  45.0000000000000]
[ 45.0000000000000  45.0000000000000  1065.00000000000 -1155.00000000000]
sage: M=matrix(QQ,[[-1155,45,45,1065], [1065,-1155,45,45], [45,1065,-1155,45], [45,45,1065,-1155]]) 
sage: M.echelon_form()       
[ 1  0  0 -1]
[ 0  1  0 -1]
[ 0  0  1 -1]
[ 0  0  0  0]
sage: v=vector([1,1,1,1])
sage: (v/v.norm()).n()
(0.500000000000000, 0.500000000000000, 0.500000000000000, 0.500000000000000)
\end{lstlisting}
    \end{answer}
  \item \cite{BryanLeise}
    给出这个网页网络的重要度排名。
    \begin{center}
      \includegraphics{jc/mp/ch5.10}
    \end{center}
    \begin{exparts}
      \item 取 $\alpha=0.85$。
      \item 取 $\alpha=0.95$。
      \item 注意到虽然所有其他网页都链接到 $p_3$，
        因而它看似重要，但它却不是排名最高的网页。
        哪个网页排名最高？
        请解释。
    \end{exparts}
    \begin{answer}
      我们有下式。
      \begin{equation*}
          H=\begin{mat}
               0     &0    &1    &1/2   \\
               1/3   &0    &0    &0   \\
               1/3   &1/2  &0    &1/2 \\
               1/3   &1/2  &0    &0
          \end{mat}
      \end{equation*}
      \begin{exparts}
        \item 下面的 \Sage{} 会话给出了答案。
\begin{lstlisting}
sage: H=matrix(QQ,[[0,0,1,1/2], [1/3,0,0,0], [1/3,1/2,0,1/2], [1/3,1/2,0,0]])
sage: S=matrix(QQ,[[1/4,1/4,1/4,1/4], [1/4,1/4,1/4,1/4], [1/4,1/4,1/4,1/4], [1/4,1/4,1/4,1/4]])
sage: I=matrix(QQ,[[1,0,0,0], [0,1,0,0], [0,0,1,0], [0,0,0,1]])
sage: alpha=0.85
sage: G=alpha*H+(1-alpha)*S
sage: N=G-I
sage: 1200*N
[-1155.00000000000  45.0000000000000  1065.00000000000  555.000000000000]
[ 385.000000000000 -1155.00000000000  45.0000000000000  45.0000000000000]
[ 385.000000000000  555.000000000000 -1155.00000000000  555.000000000000]
[ 385.000000000000  555.000000000000  45.0000000000000 -1155.00000000000]
sage: M=matrix(QQ,[[-1155,45,1065,555], [385,-1155,45,45], [385,555,-1155,555],
....:      [385,555,45,-1155]])
sage: M.echelon_form()
[            1             0             0 -106613/58520]
[            0             1             0        -40/57]
[            0             0             1        -57/40]
[            0             0             0             0]
sage: v=vector([106613/58520,40/57,57/40,1])
sage: (v/v.norm()).n()
(0.696483066294572, 0.268280959381099, 0.544778023143244, 0.382300367118066)
\end{lstlisting}
        \item 继续该 \Sage{} 会话得到下式。
\begin{lstlisting}
 sage: alpha=0.95
sage: G=alpha*H+(1-alpha)*S
sage: N=G-I
sage: 1200*N
[-1185.00000000000  15.0000000000000  1155.00000000000  585.000000000000]
[ 395.000000000000 -1185.00000000000  15.0000000000000  15.0000000000000]
[ 395.000000000000  585.000000000000 -1185.00000000000  585.000000000000]
[ 395.000000000000  585.000000000000  15.0000000000000 -1185.00000000000]
sage: M=matrix(QQ,[[-1185,15,1155,585], [395,-1185,15,15], [395,585,-1185,585],
....:      [395,585,15,-1185]])
sage: M.echelon_form()
[             1              0              0 -361677/186440]
[             0              1              0         -40/59]
[             0              0              1         -59/40]
[             0              0              0              0]
sage: v=vector([361677/186440,40/59,59/40,1])
sage: (v/v.norm()).n()
(0.713196892748114, 0.249250262646952, 0.542275102671275, 0.367644137404254)
\end{lstlisting}
       \item 网页 $p_3$ 是重要的，但它的重要度只传递给
         一个网页，即 $p_1$。
         于是该网页得到很大的提升。
      \end{exparts}    
    \end{answer}
\end{exercises}
\index{page ranking|)}
\endinput
